\section{Discussion}
\label{sec:discussion}

\paragraph{What the negative results taught us.} This project's design emerged from killing naive mechanisms, and we consider the negative results as load-bearing as the positive ones. (i) \emph{Priority ordering without admission control is useless under saturation} (Section~\ref{sec:exp-fleet}) --- a work-conservation result that redirects scheduler design from ordering to shedding. (ii) \emph{Appearance embeddings fail as event memory}, offline (cross-camera label agreement 0.55--0.59 vs.\ 0.501 chance) and in-system (Table~\ref{tab:tasti}) --- memory records must carry event semantics. (iii) \emph{VLMs are near-useless as online anomaly scorers} ($\Delta$AUC $\le +0.0009$ over tier-1, alarmist with more context; Appendix~\ref{app:e0}) but \emph{effective as gated narrators} (95.5\% grounded; Table~\ref{tab:grounded}) --- the same model class, two roles, opposite verdicts; the role assignment is the design. (iv) \emph{Per-window causal adaptation captures 0\% of oracle headroom} (Appendix~\ref{app:e0}), which is why the fleet arbiter is a slow controller over write rates rather than a per-window predictor.

\paragraph{Threats to validity.} (i) VLM-direct matches \sys's balanced accuracy on video-scoped queries (Section~\ref{sec:exp-baselines}); our differentiation is economic (7.4 GPU-s/query forever vs.\ $\approx$0) and structural (no corpus queries, no memory) --- readers should weigh accuracy parity against operating cost. (ii) Tier-1 salience is replayed from causally-computed score files, not recomputed live; its measured cost is charged, but a fully live tier-1 integration remains engineering. (iii) The fleet study is a virtual-time simulator with real GPU work, not a deployment; absolute fleet numbers depend on the acceleration and event-rate assumptions we document. (iv) \smb v0's known limits: no SHT normal videos (negation there uses taxonomy mismatch), video-level XD category attribution, query-after-ingest only. (v) The generality arm uses token-F1 in place of RVS's GPT-judge (unreachable offline); absolute levels are low and only relative comparison is meaningful. (vi) Negation FAR improves mechanically as memory shrinks; tight-budget accuracy numbers carry this caveat.

\paragraph{Open problems the experiments exposed.} Two of our pre-registered targets were missed, and both are research problems, not engineering debt. (1) \emph{Sub-event temporal precision}: tier-1 score plateaus saturate, and refined spans remain $\sim$3.5$\times$ ground-truth length (tIoU 0.286 vs.\ target 0.40). Finer temporal grounding likely needs a different signal (motion boundaries, narrated sub-spans) than detector plateaus. (2) \emph{The confirmation/abstention frontier}: verification fixes hallucinated confirmations (FAR 0.44$\to$0.08) but destroys genuine confirmation (0.95$\to$0.50--0.60) because narrations are generic while queries use category vocabulary; closing this gap (better write-time specificity, calibrated verifier confidence) is a real direction.

\paragraph{Toward hybrid memory.} The generality arm shows caption memory is lossy for fine-grained procedural QA, where full-fidelity KV retrieval wins accuracy (0.212 vs.\ 0.130 token-F1) at 34$\times$ the storage. A principled architecture follows: text/tag records as the semantic tier (cheap, queryable, fleet-wide), with an optional per-stream KV tier~\cite{rekv,livevlm} for detail, both under one budget. The eviction and admission machinery this paper builds applies to both tiers unchanged.

\paragraph{Generalizability of the economics.} Our absolute numbers are V100-era; the \emph{ratios} (event-driven vs.\ time-driven write cost, MB vs.\ GB per stream-day, ordering vs.\ admission control) are architectural and should transfer to newer hardware, since all systems compared share the same hardware in our evaluation.
